\textbf{7.\ \textsc{Electrons in magnetic field} (9 points)} --- \textit{Kaarel Hänni, Jaan Kalda.}

In what follows we consider two electrons (of mass $m$ and charge $-e$) moving
in homogeneous magnetic field of flux density $B$ in such a way that the
distance between the two electrons remains always constant; different subtasks
explore different possibilities.

\textbf{i)} \textit{(2 points)} The distance between the electrons is so large
that their electrostatic repulsion can be neglected. At a certain moment, the
angle between the velocity vectors of the electrons is $\alpha \neq 0$, and one
of them moves in the direction of the other with speed $v$. Sketch the
trajectories of the both electrons. What is the speed of the other electron?

\textbf{ii)} \textit{(1 point)}The electrostatic repulsion is still to be
neglected. Now, the trajectories of the two electrons intersect, and at a
certain moment, the velocity of one of them is $\vec{v}$. What can be said
about the velocity of the other electron at the same moment? Sketch the both
trajectories.

\textbf{iii)} \textit{(2 points)} What is the distance $l$ between the
electrons if both of them perform a periodic motion of period
$\frac{6\pi m}{Be}$ while moving with a constant speed?

\textbf{iv)} \textit{(2 points)}The figure below depicts the trajectory of one
of the electrons for a certain set of initial conditions. Where is the other
electron at the moment when the first electron is at the position of the small
circle? What is the period of this motion?

\begin{center}\includegraphics[width=0.44\linewidth]{figures/NBPhO_2022_Q7_fig1.png}\end{center}

\textbf{v)} \textit{(2 points)}The figure below depicts the trajectory of one
of the electrons for a certain set of initial conditions. What is the speed of
the other electron at the moment when the first electron is at the position of
the small circle?

\begin{center}\includegraphics[width=0.43\linewidth]{figures/NBPhO_2022_Q7_fig2.png}\end{center}
